\documentclass[conference]{IEEEtran}
\IEEEoverridecommandlockouts
% The preceding line is only needed to identify funding in the first footnote. If that is unneeded, please comment it out.
\usepackage{amsmath,amssymb,amsfonts}
\usepackage{algorithmic}
\usepackage{graphicx}
\usepackage{textcomp}
\usepackage{xcolor}
\usepackage{twemojis} 
\usepackage[numbers]{natbib}

\usepackage{booktabs} % For professional looking tables
\usepackage{hyperref} % For clickable links and citations
\usepackage{enumitem}   % Allows fine-tuning of lists inside tables to save space

\def\BibTeX{{\rm B\kern-.05em{\sc i\kern-.025em b}\kern-.08em
    T\kern-.1667em\lower.7ex\hbox{E}\kern-.125emX}}
\begin{document}

\title{Sentiment Analysis of Code-Mixed Singlish (Sinhala-English) Reviews with Emoji Fusion NLP Approach\\
}

\author{\IEEEauthorblockN{K. M. Sajana Karunanayake - 268245C}
      \IEEEauthorblockA{\textit{University of Moratuwa} \\
            sajanak.26@cse.mrt.ac.lk | karunanayakekmsn.26@uom.lk}
}

\maketitle

\section{Introduction}
% % TODO: Motivate the problem -- rise of code-mixed, emoji-rich user-generated content in Sri Lankan e-commerce/app ecosystems; why sentiment analysis matters for businesses and platforms; why Sinhala/Singlish + emoji is an underserved combination.
In today's Srilanka Almost every E-Commerce platforms and social media have
seen a widespread adoption of ``Singlish'' a code-mixed combination pf Sinhala
+ English rendered in Roman Script \cite{uthpala2024sinhala}. Standard
sentiment analysis tools designed for monolingual formal text perform very
poorly on Singlish due to fluid intra-sentential Code-switching, informal
spelling variations and High morphological complexity
\cite{chakravarthi2020corpus}. Recent natural language processing (NLP)
research has increasingly explored multilingual pre trained transformer encoder
for example \textbf{XLM-RoBERTa} for low resource Indian and South Asian
languages \cite{devlin2019bert,desilva2019survey}, User feedback online relies
heavily on non verbal, visual affective markers specifically emojis. Emojis
condense sentiment, emotional tone And sarcasm across language barriers
sentiment, emotional tone, and sarcasm across language barriers
\cite{barbieri2017are,eisner2016emoji2vec}. Recent multimodal frameworks
demonstrate that fusing textual representations with dense emoji embeddings
significantly enchances classification accuracy in codemixed scenarios
\cite{singh2024predicting}. In this project we check multimodel sentiment
classification for Singlish product and application reviews by fusing
contextual transformer representations with pre trained emoji embeddings.

\section{Research problem}

Now existing researches on the Singlish Sentiment analysis routinely discard
emojis, emoticons, non alphanumeric symbols during text preprocessing,
pinpointing them as irrecoverable noise or rather ignored
\cite{uthpala2024sinhala,desilva2019survey}. This raises Critical
\textbf{research Gaps}:

\begin{enumerate}
      \item \textbf{Loss of essential Affective Signals:}
            Emojis Carries vital emotional polarity and sarcasm in consumer reviews. For example Brief
            Singlish phrase paired with a fire \twemoji{fire} or with crying emoji \twemoji{loudly crying face}
            Which is completely lost when the emojis are removed.
      \item \textbf{Ommision of deep feature Fusion:}
            Current Singlish corpora rely mosttly exclusively in classical machine learning baselines
            (Logistic Regression, SVM) \cite{uthpala2024sinhala} or unadapted multilingual BERT without without
            cross-modal feature fusion \cite{chakravarthi2020corpus}
      \item \textbf{Data sourfce \& Domain Limitations:} Existing Singlish Datasets are restricted to Youtube
            comments \cite{uthpala2024sinhala} leaving e-commerce product's feedbacks Like Daraz.lk, Social media
            marketplaces Like Facebook marketplace, Mobile app store reviews?.
\end{enumerate}

Solving this problem is essential to enable accurate and automated market
intelligence, product reputations tracking, customer experience analytics in
Sri Lanka's digital economy.

\section{Research Objectives}
% % TODO: List 3-5 concrete objectives, e.g.:
% % 1. Construct and publicly release an annotated Singlish-emoji review corpus.
% % 2. Benchmark multilingual transformer models (XLM-R, etc.) on this corpus.
% % 3. Design and evaluate emoji-text fusion strategies.
% % 4. Analyze error patterns specific to code-mixing and emoji ambiguity.
Primary Objective if this project is to develop, evaluate and release a
benchmark emoji-fused sentiment analysis framework for CodeMixed singlish
consumer reviews. The specific research goals are:

\begin{itemize}
      \item \textbf{01 (Dataset Construction):} Harvest and annotate a new, gold standard Singlish review dataset
            from e-commerce sites, mobile app stores, facebook marketplace reviews that explicitly retains all
            inline Unicode emojis.
      \item \textbf{02 (Representational learning):} Extract contextual textual features by using fine tuned multilingual
            encoders (XLM-RoBERTa) and map them alongside 300-dimential emoji vector derived from \texttt{emoji2vec} and
            Unicode description specs
      \item \textbf{03 (Fusion Architecture):} Design and implement a Cross Modal Emoji Text Fusion Layer that eveluates
            Gated feature Fusion and cross attention modules  to dynamically combine textual and visual affective signals.
      \item \textbf{04 (Empirical Evaluation):} Benchmarking the fused model against text only transformer baseline and classical
            ML classifiers across Accuracy, Macro F1 Score and per class sentiment metrics (Positive, Negative Neutral).
\end{itemize}

\section{Proposed method}
% % TODO: Cover (a) Sinhala/Singlish NLP and code-mixed sentiment analysis generally, (b) emoji sentiment/embedding research, (c) multilingual transformer fine-tuning (XLM-R, mBERT), (d) existing low-resource review datasets, and identify the specific gap this thesis fills.
\textbf{Main Concept:} The proposed solution uses a two stream deep learning network. The \textit{Text Steam} Passes raw Singlish token
sequences into a fine tuned XLM-RoBERTa encoder. The \textit{Emoji Stream} extracts 300-dimensional continous representations for detected
emojis usaing pre trained \texttt{emoji2vec} \cite{eisner2016emoji2vec} and multi label emoji embeddings \cite{singh2024predicting}. A Gated fusion model
elarns dynamic attention weights to merge text and emoji representations into a unified feateure vector prior to softmax entiment classifications.

% \section{Methodology}
% \subsection{Data Collection}
% % TODO: Scraping approach for Daraz.lk product reviews and Sri Lankan mobile app store reviews (Google Play/App Store); note PickMe has no public review feature and is excluded; ethical/ToS considerations; target corpus size.
% \subsection{Annotation}
% % TODO: Annotation scheme (e.g., positive/negative/neutral, or fine-grained), number of annotators, inter-annotator agreement metric (e.g., Cohen's/Fleiss' kappa).
% \subsection{Model Development}
% % TODO: Baselines (lexicon-based, classical ML); fine-tuning XLM-R or similar multilingual transformer; emoji representation (emoji2vec, learned embeddings) and fusion strategy (early/late fusion, attention-based).
% \subsection{Evaluation}
% % TODO: Metrics (accuracy, macro-F1), ablations (text-only vs. emoji-fused), error analysis plan.

% \section{Expected Contributions}
% % TODO: (1) Novel public dataset; (2) benchmark results for code-mixed Sinhala sentiment analysis; (3) emoji fusion methodology generalizable to other low-resource code-mixed settings.

\clearpage
\nocite{*}
\bibliographystyle{IEEEtranN} % or 'unsrtnat'
\bibliography{references}
\end{document}